%%% ============================================================
%%% The 70/71/72/73 HARMONY Ladder:
%%%   Three Independent Algebraic Constructions and One Corollary
%%%   on Z/10Z
%%%
%%% Brayden Ross Sanders, M. Gish
%%% 7Site LLC, Hot Springs AR, 2026
%%% DOI: 10.5281/zenodo.18852047
%%%
%%% Target venue:  Journal of Combinatorial Theory, Series A (JCT-A)
%%% Source:        FORMULAS_AND_TABLES.md §J D97;
%%%                Gen13/targets/foundations/tables/harmony_ladder.py
%%% MSC:           05A05, 05E18, 11T22, 16Y99, 20N02
%%%
%%% Revision history (2026-05-07): defensive-exposition pass per
%%% J22_JCT-A_FreshEyes_REBUTTAL.md.
%%%   (a) Title sharpened: "Three Independent Constructions and One
%%%       Corollary" replaces "Four Independent Constructions"; the
%%%       70-rung is positioned explicitly as a corollary rather than
%%%       an independent construction (§5; the 71-rung's three roles
%%%       remain the genuinely-independent triple).
%%%   (b) Embedded explicit sympy verification snippet (§A) for the
%%%       discriminant claim disc(x^4+4x^3-x^2+2x-2) = -40896
%%%       = -2^6 * 3^2 * 71. Pre-empts future referee factoring
%%%       errors (the original referee's claimed factorization
%%%       2^7 * 3 * 7 * 19 = 51072 is arithmetically wrong; the
%%%       magnitudes don't match -40896, off by 25%).
%%%   (c) Verification scripts in manuscript/verification/ replaced
%%%       per save plan: the so10/Lie-algebra scripts that were in
%%%       the folder by mistake have been moved aside, and the
%%%       canonical HARMONY-ladder scripts (tsml_harmony_count.py,
%%%       tsml_submagma_9x9.py, tsml_bhml_disagreement.py,
%%%       bhml_8_ym_det.py, plus harmony_ladder_disc_check.py for
%%%       the discriminant) are bundled.
%%% ============================================================

\documentclass[11pt,reqno]{amsart}

\usepackage{amsmath,amssymb,amsthm}
\usepackage[margin=1in]{geometry}
\usepackage{hyperref}
\usepackage{enumitem}
\usepackage{array}
\usepackage{booktabs}
\usepackage{listings}
\usepackage{xcolor}

\lstdefinestyle{plainpy}{
  basicstyle=\ttfamily\footnotesize,
  language=Python,
  showstringspaces=false,
  breaklines=true,
  frame=single,
  numbers=none,
}

\newtheorem{theorem}{Theorem}[section]
\newtheorem{lemma}[theorem]{Lemma}
\newtheorem{proposition}[theorem]{Proposition}
\newtheorem{corollary}[theorem]{Corollary}

\theoremstyle{definition}
\newtheorem{definition}[theorem]{Definition}
\newtheorem{remark}[theorem]{Remark}
\newtheorem{example}[theorem]{Example}

\newcommand{\Z}{\mathbb{Z}}
\newcommand{\R}{\mathbb{R}}
\newcommand{\F}{\mathbb{F}}
\newcommand{\TSML}{\mathrm{TSML}}
\newcommand{\BHML}{\mathrm{BHML}}
\DeclareMathOperator{\HARM}{HARM}
\DeclareMathOperator{\disc}{disc}

\title[The 70/71/72/73 HARMONY Ladder]
  {The 70/71/72/73 HARMONY Ladder:\\
   Three Independent Algebraic Constructions and One Corollary
   on $\Z/10\Z$}

\author{Brayden R.\ Sanders \and M.\ Gish}
\address{7Site LLC, Hot Springs, Arkansas, USA}
\email{brayden@7site.co}

\author{Brayden R.\ Sanders \and M.\ Gish}
\address{Independent Researcher, Hot Springs, Arkansas, USA}
\email{monica.gish1992@gmail.com}

\date{2026}

\keywords{finite magma, composition table, sub-magma, integer
  invariant, structural clustering, Galois prime, Yang--Mills core,
  $E_6$ root system}

\subjclass[2020]{05A05, 05E18, 11T22, 16Y99, 20N02}

\begin{document}

\begin{abstract}
We record an integer-invariant ladder on the canonical composition
lattice over $\Z/10\Z$ whose values cluster at $\{70, 71, 72, 73\}$.
The three structurally-independent rungs are HARMONY-cell counts of
distinct sub-magmas: the full $10 \times 10$ TSML table $T$ has
$\HARM(T) = 73$; the VOID-stripped $9 \times 9$ sub-magma has
$\HARM(T_{\{1, \ldots, 9\}}) = 71$; and the canonical companion table
$B$ on the Yang--Mills $8 \times 8$ sub-matrix has $\det = 70 = \binom{8}{4}$.
The $72$-rung is a corollary of the $73$-rung (the $(7,7)$ self-cell
removed). The integer $71$ enters the ladder in three
independently-verified structural roles simultaneously: as the HARMONY
count of the VOID-stripped $9 \times 9$ sub-magma; as the
cell-disagreement count between the two canonical $\Z/10\Z$ composition
tables; and as the unique odd prime $>3$ appearing in the polynomial
discriminant $\Delta_f = -2^6 \cdot 3^2 \cdot 71 = -40896$ (equivalently
in the field discriminant $d_K = -2^4 \cdot 3^2 \cdot 71 = -10224$) of
the quartic number field LMFDB $4.2.10224.1$ that governs the
closed-form runtime attractor on the same substrate.
We verify each construction at integer/machine precision (sympy
\texttt{discriminant} and \texttt{factorint}; numpy direct count) and
embed the verification snippets inline. The results are stated as
Tier-B forced consequences of the canonical table structure; no
axiom-level forcing is required.
\end{abstract}

\maketitle

\section{Introduction}
\label{sec:intro}

Combinatorial substrates over small finite cyclic groups occasionally
exhibit clusters of integer invariants whose values agree, or nearly
agree, across structurally distinct constructions. When the
clustering is tight enough that the same prime, or four consecutive
integers, recur across independent algebraic objects on the same
substrate, the clustering is itself a structural invariant of the
substrate.

This paper records one such clustering on the canonical composition
lattice over $\Z/10\Z$. The substrate carries two related
$10 \times 10$ composition tables, the \emph{TSML} table $T$ and the
\emph{BHML} table $B$, both fixed by the canonical CL forcing axioms
recorded in \cite{Sanders2026CLAxioms}. On this substrate three
\emph{structurally-independent} integer invariants and one
\emph{inclusion-exclusion corollary} arise from algebraic constructions:
\begin{itemize}[leftmargin=*]
\item \emph{Independent rung A.} $\HARM(T) = 73$ (HARMONY-cell count
  of the full TSML table);
\item \emph{Corollary rung.} $\HARM(T) - 1 = 72$ (drop the $(7,7)$
  self-cell apex of the rung-A count; equals the number of positive
  roots of $E_6$ as a numerical coincidence — see
  Remark~\ref{rem:e6});
\item \emph{Independent rung B.} $\HARM(T_{\{1,\ldots,9\}}) = 71$
  (HARMONY count of the VOID-stripped $9 \times 9$ sub-magma);
  equivalently $|T \ominus B| = 71$ (the cell-disagreement count
  between the two canonical $\Z/10\Z$ composition tables);
  equivalently the unique odd prime $> 3$ in the polynomial
  discriminant $\Delta_f = -2^{6} \cdot 3^{2} \cdot 71 = -40896$ of
  the quartic $f(x) = x^4 + 4 x^3 - x^2 + 2 x - 2$ defining the
  number field LMFDB $4.2.10224.1$ (whose field discriminant is
  $d_K = -2^{4} \cdot 3^{2} \cdot 71 = -10224$, i.e.\ the
  ``$10224$'' in the LMFDB identifier);
\item \emph{Independent rung C.} $\det(B_{\{1,2,3,4,5,6,8,9\}}) = 70 = \binom{8}{4}$
  (the Yang--Mills core sub-determinant).
\end{itemize}

The four invariants are pairwise distinct integers in
$\{70, 71, 72, 73\}$, three of them obtained from genuinely independent
constructions and one from inclusion-exclusion on the largest. The
clustering at four consecutive integers is the content of the ladder.
The triple-coincidence at $71$ (sub-magma HARMONY count $=$
lens-disagreement count $=$ Galois prime) is the sharpest single point
of the ladder; we verify all three constituent claims independently
(§\ref{sec:71}) and pre-empt potential factoring errors with an
embedded sympy verification snippet for the discriminant claim
(§\ref{subsec:71-galois}, §\ref{sec:verification-snippet}).

\paragraph{Statement of the main theorem.}

Define the ladder $\mathcal{L} = (73, 72, 71, 70)$ as the
$4$-tuple of invariants arising from the constructions enumerated
above. Then:

\begin{theorem}[Ladder verification]
\label{thm:ladder}
Each of the four entries of $\mathcal{L}$ equals the value computed
by direct integer enumeration on its construction. The ladder is
exact at machine precision, with each entry verified by a sympy /
numpy snippet (§\ref{sec:verification-snippet}). Three of the four
constructions (rungs A, B, C) are structurally independent in the
sense of Remark~\ref{rem:independence} below; the $72$-rung is an
inclusion-exclusion corollary of the $73$-rung.
\end{theorem}

\begin{remark}[Independence of the three structurally-independent
rungs]
\label{rem:independence}
Rungs A ($\HARM(T) = 73$), B ($\HARM(T_{\{1,\ldots,9\}}) = 71$, with
its three structurally-distinct constituent constructions), and C
($\det(B_{\mathrm{YM}}) = 70$) are independent in the sense that none
is deducible from any other by elementary cell-counting on the same
sub-matrix; each requires its own combinatorial or linear-algebraic
input. The $72$-rung is the inclusion-exclusion consequence of the
$73$-rung after excising the single cell $(7, 7)$. Within rung B,
the three constituent constructions are themselves independent (a
HARMONY-cell count on $T$ restricted to a $9 \times 9$ sub-magma; a
bitwise comparison of $T$ against $B$; and an algebraic-number-theory
discriminant computation) and converge on the same prime $71$.
\end{remark}

The clustering at $\{70, 71, 72, 73\}$ is therefore a non-trivial
structural invariant: three independent constructions plus one
inclusion-exclusion corollary, with the prime $71$ entering through
three further structurally-distinct mechanisms inside rung B.

\paragraph{Outline.}

Section~\ref{sec:setup} fixes the substrate, the two composition
tables $T$ and $B$, and the HARMONY operator. Section~\ref{sec:73}
establishes the $73$-rung. Section~\ref{sec:72} establishes the
$72$-rung-corollary and records the $E_6$ coincidence.
Section~\ref{sec:71} establishes the $71$-rung in three structurally
distinct ways and proves their numerical agreement.
Section~\ref{sec:70} establishes the $70$-rung and records the
$\binom{8}{4}$ coincidence. Section~\ref{sec:companion} records two
companion HARMONY counts at $\{28, 36, 44\}$ that lie on a parallel
ladder structure (each appearing in two structural roles, not four).
Section~\ref{sec:reading} records the structural reading of the
ladder as a descent. Section~\ref{sec:verification-snippet} embeds
the sympy/numpy verification snippets; the discriminant snippet of
§\ref{subsec:disc-snippet} is the one we recommend any reader run
first.

\paragraph{Lens and substrate.}
\label{sec:lens-subs}

We work on $\Z/10\Z$ with the specific 10-operator labels recorded
above and the canonical $(T, B) = (\TSML, \BHML)$ table pair fixed by
the CL forcing axioms~\cite{Sanders2026CLAxioms}. These choices are
\emph{not derived from first principles} in this paper; they reflect a
structural reading of the substrate motivated by the source program's
10-operator decomposition and the observed dynamics. The theorems
below are theorems on this specific structure; analogous theorems
would hold on other small finite commutative non-associative
substrates with appropriately chosen tables. The framing follows the
Drápal--Wanless (2021)~\cite{DrapalWanless2021} line of work on small
finite commutative non-associative structures with structural
invariants.

\paragraph{PROVEN / COMPUTED / STRUCTURAL RHYME / OPEN.}

\textbf{PROVEN:} the four-rung ladder $(73, 72, 71, 70)$ on the canonical
$(T, B)$ tables; in particular, the triple-coincidence at the
71-rung (sub-magma HARMONY count $=$ lens-disagreement count $=$ unique
odd Galois prime $>3$ in the polynomial discriminant $\Delta_f$ of the
quartic defining LMFDB $4.2.10224.1$). \textbf{COMPUTED:} all four rungs
verified at integer/machine precision by the embedded sympy/numpy
snippets of §\ref{sec:verification-snippet}; the discriminant
factorization $-40896 = -2^6 \cdot 3^2 \cdot 71$ is confirmed by
\texttt{sympy.discriminant} and \texttt{sympy.factorint} in
§\ref{subsec:disc-snippet}. \textbf{STRUCTURAL RHYME:} the
$E_6^+ = 72$ and $\binom{8}{4} = 70$ identifications are recorded as
numerical coincidences with adjacent algebra/representation theory;
no derivation is claimed. \textbf{OPEN:} a structural explanation for
the convergence of three independent constructions on the same prime
$71$ at the rung-B level (i.e., a single structural object yielding
all three constructions as projections); whether the ladder extends
beyond rungs A--C to a longer descent.

\paragraph{Lens scope.}

The TSML table admits two principled lens-symmetrization choices,
$T_{\mathrm{RAW}}$ (literal, non-commutative) and $T_{\mathrm{SYM}}$
(upper-triangle symmetrized, commutative). All four ladder rungs in
this paper are \emph{lens-invariant} on both lenses:
$\HARM(T_{\mathrm{RAW}}) = \HARM(T_{\mathrm{SYM}}) = 73$, and the
sub-magma counts at $9 \times 9$ are likewise unchanged. The wobble
phenomenon (prime $11$ in the characteristic polynomial coefficients)
distinguishing RAW from SYM is the subject of a separate
companion \cite{Sanders2026Wobble}; it does not affect the ladder.

\section{Setup}
\label{sec:setup}

The substrate is $\Z/10\Z = \{0, 1, \ldots, 9\}$ with the operator
labels recorded in \cite{Sanders2026CLAxioms}: $\mathbf{V} = 0$
(VOID), $\mathbf{L} = 1$, $\mathbf{C} = 2$, $\mathbf{P} = 3$,
$\mathbf{Co} = 4$, $\mathbf{Ba} = 5$, $\mathbf{Ch} = 6$, $\mathbf{H}
= 7$ (HARMONY), $\mathbf{Br} = 8$, $\mathbf{R} = 9$.

The TSML composition table $T \in \Z^{10 \times 10}$ takes values in
$\Z/10\Z$, with $T(i, j)$ the composition of operators $i$ and $j$.
By the CL forcing axioms (\cite{Sanders2026CLAxioms}, Axioms A1--A9),
$T$ is uniquely determined; we use the upper-triangle symmetrized
form $T_{\mathrm{SYM}}$ throughout this paper unless noted otherwise.
Define
\[
  \HARM(T) \;=\; |\{(i, j) : T(i, j) = 7\}|.
\]

The BHML composition table $B$ is the companion table prescribed by
the same axiom system at the curvature lens; its construction is
recalled in \cite{Sanders2026LensInvariance}. We will use $B$ in two
of the four ladder constructions.

For any subset $S \subseteq \{0, \ldots, 9\}$ that is closed under
$T$ (a sub-magma), define
\[
  T_S \;=\; T |_{S \times S},\qquad
  \HARM(T_S) \;=\; |\{(i, j) \in S \times S : T(i, j) = 7\}|.
\]
The same definitions apply to $B$. We do not require $S$ to be a
sub-magma of $B$ for the determinant constructions on $B$; we only
require $S$ to specify a square sub-matrix.

\section{The 73-rung: full HARMONY count}
\label{sec:73}

\begin{theorem}[$73$-rung]
\label{thm:73}
$\HARM(T) = 73$.
\end{theorem}

\begin{proof}
The HARMONY-cell decomposition is given by the three disjoint
exception classes proved in \cite{Sanders2026CLAxioms} Theorem D10:
\begin{itemize}[leftmargin=*]
\item the VOID-row $V_0 = \{(0, j) : j \neq 7\}$ contributes $9$ non-HARMONY cells;
\item the VOID-column $V_1 = \{(i, 0) : i \in \{1,\ldots,6,8,9\}\}$ contributes $8$ further non-HARMONY cells (cell $(0,0)$ already in $V_0$, cell $(7,0)$ is HARMONY);
\item the ECHO pairs $\{(1,2),(2,1),(2,4),(4,2),(2,9),(9,2),(3,9),(9,3),(4,8),(8,4)\}$ contribute $10$ non-HARMONY cells, and these are disjoint from $V_0 \cup V_1$.
\end{itemize}
Total non-HARMONY: $9 + 8 + 10 = 27$. Therefore
$\HARM(T) = 100 - 27 = 73$.
\end{proof}

\section{The 72-rung: drop the apex}
\label{sec:72}

\begin{theorem}[$72$-rung]
\label{thm:72}
$\HARM(T) - 1 = 72$, where the subtracted cell is $T(7, 7) = 7$
(the HARMONY self-cell apex).
\end{theorem}

\begin{proof}
Direct: $T(7, 7) = 7$ by Axiom A7 (HARMONY is absorbing in row $7$
column $7$). Removing this cell from the count yields $73 - 1 = 72$.
\end{proof}

\begin{remark}[$E_6$ coincidence]
\label{rem:e6}
The simple Lie algebra $E_6$ has exactly $72$ positive roots
\cite{Bourbaki1968,FultonHarris1991}. We record the coincidence
without claiming a derivation. In particular, the embedding of the
TSML composition lattice into a root-system framework is not yet
established; the coincidence enters the ladder as a numerical
fact, not as a structural identification.
\end{remark}

\section{The 71-rung: three structural roles}
\label{sec:71}

The integer $71$ enters the ladder via three independent constructions
that all yield the same value. We state and prove each.

\begin{theorem}[$71$-rung, sub-magma form]
\label{thm:71-submagma}
$\HARM(T_{\{1,\ldots,9\}}) = 71$, where
$T_{\{1,\ldots,9\}}$ is the restriction of $T$ to the
VOID-stripped $9 \times 9$ sub-magma.
\end{theorem}

\begin{proof}
Restricting $T$ to $\{1, \ldots, 9\} \times \{1, \ldots, 9\}$
removes the VOID-row $V_0$ and the VOID-column $V_1$ from the
non-HARMONY exception classes, but the ECHO pairs (which lie
entirely in $\{1, \ldots, 9\} \times \{1, \ldots, 9\}$) remain.
The reduced count of non-HARMONY cells in the $9 \times 9$ matrix
is therefore exactly $|\mathrm{ECHO}| = 10$ minus zero (none of
the ECHO pairs touch row $0$ or column $0$). The total cell count
is $9 \times 9 = 81$, and the HARMONY count is
$81 - 10 = 71$.
\end{proof}

\begin{theorem}[$71$-rung, lens-disagreement form]
\label{thm:71-disagree}
$|\{(i, j) : T(i, j) \neq B(i, j)\}| = 71$.
\end{theorem}

\begin{proof}
Direct cell-by-cell comparison of the two canonical tables; the
$71$-cell disagreement count is verified at machine precision in
the script \texttt{tsml\_bhml\_disagreement.py}; see also
\cite{Sanders2026LensInvariance}.
\end{proof}

\begin{theorem}[$71$-rung, Galois form]
\label{thm:71-galois}
\label{subsec:71-galois}
$71$ is the unique odd prime $> 3$ in the polynomial discriminant
$\Delta_f = \disc(f) = -40896 = -2^{6} \cdot 3^{2} \cdot 71$ of the
quartic $f(x) = x^4 + 4 x^3 - x^2 + 2 x - 2 \in \Z[x]$. The same
prime $71$ appears in the field discriminant $d_K = -10224 = -2^4
\cdot 3^2 \cdot 71$ of $K = \Q[x]/(f)$ (with $|\Delta_f|/|d_K| = 4$,
so $[\mathcal{O}_K : \Z[\xi^*]] = 2$); the magnitude $|d_K| = 10224$
is the third coordinate of the LMFDB identifier $4.2.10224.1$
\cite{Sanders2026Attractor}. The polynomial $f$ is identified in
\cite{Sanders2026Attractor} as the minimal polynomial of the
coordinate ratio $\xi^* = r/\beta$ of the closed-form $\TSML+\BHML$-mix
runtime attractor at mixing weight $\alpha = 1/2$, and
$\mathrm{Gal}(f/\Q) = D_4$.
\end{theorem}

\begin{proof}
The discriminant is computed from the minimal polynomial of the
attractor coordinate ratio $H/\mathrm{Br} = 1 + \sqrt{3}$ together
with the second algebraic ratio identified by the PSLQ analysis in
\cite{Sanders2026Attractor}; the LMFDB identification $4.2.10224.1$
is confirmed by direct comparison. The factorization
$-40896 = -2^6 \cdot 3^2 \cdot 71$ is verified at integer precision
by the sympy snippet below; an alternative factorization claim of
$-40896 = -2^7 \cdot 3 \cdot 7 \cdot 19$ would equal $-51072$, which
differs from $-40896$ by approximately $25\%$ and is therefore
arithmetically wrong, not a difference of convention.

\begin{lstlisting}[style=plainpy]
import sympy
x = sympy.Symbol('x')
f = x**4 + 4*x**3 - x**2 + 2*x - 2
disc = sympy.discriminant(f, x)
print("disc =", disc)                   # -40896
print("factor:", sympy.factorint(abs(disc)))
                                         # {2: 6, 3: 2, 71: 1}
\end{lstlisting}
\end{proof}

\begin{corollary}[Triple coincidence at $71$]
\label{cor:71-triple}
The integer $71$ appears in three independently-verified structural
roles on the same substrate: as the HARMONY count of the VOID-stripped
$9 \times 9$ sub-magma (Theorem~\ref{thm:71-submagma}), as the
cell-disagreement count between the two canonical lens tables
(Theorem~\ref{thm:71-disagree}), and as the unique odd prime in the
Galois discriminant of the quartic field governing the substrate's
closed-form attractor (Theorem~\ref{thm:71-galois}).
\end{corollary}

\begin{remark}[Reading the triple coincidence]
\label{rem:71-reading}
The three constructions in Corollary~\ref{cor:71-triple} are
genuinely independent: the first is pure cell-counting on $T$, the
second is bitwise comparison of $T$ against $B$, and the third is a
discriminant computation in algebraic number theory keyed to the
runtime attractor. Each was discovered separately in the development
of the substrate; their numerical agreement at $71$ was not assumed
in any of the three constructions and is not a corollary of a
shared structural identity.
\end{remark}

\section{The 70-rung: Yang--Mills determinant}
\label{sec:70}

The fourth ladder rung is not a HARMONY count but a determinant. It
arises from the $8 \times 8$ sub-matrix of $B$ obtained by dropping
the VOID and HARMONY indices from $B$:
\[
  B_{\mathrm{YM}} \;=\; B|_{\{1,2,3,4,5,6,8,9\} \times \{1,2,3,4,5,6,8,9\}}.
\]

\begin{theorem}[$70$-rung]
\label{thm:70}
$\det(B_{\mathrm{YM}}) = 70 = \binom{8}{4}$.
\end{theorem}

\begin{proof}
The determinant is computed by direct integer expansion of the
$8 \times 8$ matrix $B_{\mathrm{YM}}$; the value $70$ is exact in
$\Z$. The combinatorial identity $\binom{8}{4} = 70$ is elementary.
The script \texttt{bhml\_8\_ym\_det.py} performs the computation at
integer precision.
\end{proof}

\begin{remark}[$\binom{8}{4}$ structural reading]
\label{rem:70-reading}
The integer $\binom{8}{4} = 70$ is the dimension of the space of
$4$-forms on $\R^8$, and equivalently the number of independent
self-dual or anti-self-dual $4$-form components on $\R^8$. We record
the coincidence between $\det(B_{\mathrm{YM}})$ and $\binom{8}{4}$
as a numerical fact; the labelling ``$\mathrm{YM}$'' for the
$\{1, 2, 3, 4, 5, 6, 8, 9\}$ sub-matrix reflects the index set's role
in the substrate's broader research program (an internal working
document on the substrate's possible bridge to Yang--Mills lattice
gauge structure, available from the authors on request) and is not
itself part of the present claim. The determinant rung sits
\emph{below} the three HARMONY-count rungs in the ladder: it lives
in a different invariant layer of the substrate, but its integer
value is the next step in the descent $73 \to 72 \to 71 \to 70$.
\end{remark}

\section{Companion HARMONY counts: a second ladder}
\label{sec:companion}

In addition to the four-rung ladder $(73, 72, 71, 70)$, the substrate
carries three further integers $(28, 36, 44)$ that each appear at
two structurally distinct roles. These are not part of the main
ladder but are recorded for completeness; their proofs are
direct counting arguments in the same spirit.

\begin{proposition}[$28$-rung companion]
\label{prop:28}
$\HARM(B) = 28$. The same integer is the dimension of the simple Lie
algebra $\mathfrak{so}(8)$.
\end{proposition}

\begin{proposition}[$36$-rung companion]
\label{prop:36}
$\HARM(T_{S_7}) = 36$, where $S_7 = \{0, 4, 5, 6, 7, 8, 9\}$ is the
size-$7$ sub-magma in the joint $T+B$ closure chain
\cite{Sanders2026FourCore}. The same integer is the BHML
$\sigma^2$-cycle-A projection count
\cite{Sanders2026LensInvariance}.
\end{proposition}

\begin{proposition}[$44$-rung companion]
\label{prop:44}
$\HARM(\mathrm{CL\_STD}) = 44$, where $\mathrm{CL\_STD}$ is the
Standard encoding table on $\Z/10\Z$ from
\cite{Sanders2026LensInvariance}, distinct from both $T$ and $B$.
The same integer is the BHML $\sigma^2$-cycle-B projection count
$28 + 11 + 5 = 44$.
\end{proposition}

The proofs are direct counting; verification scripts run in under
$1$ second each at machine precision.

\section{Structural reading and discussion}
\label{sec:reading}

The ladder $(73, 72, 71, 70)$ admits a natural reading as a four-step
descent through the substrate's HARMONY structure:

\begin{enumerate}[label=(\arabic*),leftmargin=*]
\item \textbf{$73$ — full HARMONY apex.} The complete prescribed view
  of the substrate's HARMONY-cell density.
\item \textbf{$72$ — drop the $(7,7)$ apex.} The "BEING shell" view
  with the central self-cell removed; coincides with $|E_6^+|$, the
  positive root count of $E_6$.
\item \textbf{$71$ — VOID-stripped HARMONY $=$ lens-disagreement $=$
  Galois prime.} The triple coincidence: three independent
  constructions on the substrate land on the same prime. Three
  unrelated algebras pointing at the same integer is the algebraic
  shape of a real invariant.
\item \textbf{$70$ — Yang--Mills core determinant.} The descent from
  HARMONY-cell counting into the determinant-invariant layer; the
  rung that lives one floor below but is structurally adjacent.
\end{enumerate}

The descent reflects the dominance of the HARMONY axis ($7$) on the
substrate: every nearby integer invariant lives within a few of
$\HARM(T) = 73$. The four rungs are all integers; three of them are
HARMONY counts; one of them is a determinant; the gaps between
consecutive rungs are $\{1, 1, 1\}$.

\paragraph{What becomes plain only when all four are held together.}

No single rung individually distinguishes the substrate from a
generic $10 \times 10$ table; HARMONY counts and small determinants
are common in finite-magma combinatorics. What is unusual is the
clustering at four consecutive integers from genuinely independent
constructions, together with the prime $71$ entering through three
separate mechanisms simultaneously. The Monte Carlo significance of
the full HARMONY count $\HARM(T) = 73$ alone, against $200{,}000$
random tables with the same row/column constraints, is
$Z = 21.3$, $p < 10^{-50}$ \cite{Sanders2026CLAxioms}. The remaining
three rungs reinforce this baseline: each is a non-generic invariant
of the same substrate.

\section{Embedded verification snippets}
\label{sec:verification-snippet}

The four ladder rungs are verified by short sympy / numpy snippets
that any referee can run in under two seconds with a Python REPL.
We embed the snippets directly so that no external script is required
for the load-bearing claims.

\subsection{Discriminant claim ($71$-rung, Galois form)}
\label{subsec:disc-snippet}

\begin{lstlisting}[style=plainpy]
import sympy
x = sympy.Symbol('x')
f = x**4 + 4*x**3 - x**2 + 2*x - 2
disc = sympy.discriminant(f, x)
print("disc(f) =", disc)                          # -40896
print("|disc(f)| =", abs(disc))                   # 40896
print("factor:", sympy.factorint(abs(disc)))      # {2: 6, 3: 2, 71: 1}
# Cross-check (negation of the alternate-factorization claim):
print(2**7 * 3 * 7 * 19)                          # 51072 (NOT 40896)
print(2**6 * 3**2 * 71)                           # 40896 (correct)
\end{lstlisting}

The two final lines are intentionally redundant: they let any reader
verify by direct multiplication that $-2^6 \cdot 3^2 \cdot 71 = -40896$
and that an alternative claim $-2^7 \cdot 3 \cdot 7 \cdot 19 = -51072$
disagrees with $-40896$ by approximately $25\%$ — i.e., the alternate
factorization is arithmetically wrong, not a difference of convention.

\subsection{HARMONY-count claims ($73$-rung, $71$-rung sub-magma form)}
\label{subsec:harm-snippet}

\begin{lstlisting}[style=plainpy]
import numpy as np
T = np.array([
    [0,0,0,0,0,0,0,7,0,0],
    [0,7,3,7,7,7,7,7,7,7],
    [0,3,7,7,4,7,7,7,7,9],
    [0,7,7,7,7,7,7,7,7,3],
    [0,7,4,7,7,7,7,7,8,7],
    [0,7,7,7,7,7,7,7,7,7],
    [0,7,7,7,7,7,7,7,7,7],
    [7,7,7,7,7,7,7,7,7,7],
    [0,7,7,7,8,7,7,7,7,7],
    [0,7,9,3,7,7,7,7,7,7],
])
print("HARM(T) =", int((T == 7).sum()))           # 73
T9 = T[1:, 1:]
print("HARM(T_{1..9}) =", int((T9 == 7).sum()))   # 71
\end{lstlisting}

\subsection{Lens-disagreement count ($71$-rung, lens form)}
\label{subsec:lens-snippet}

\begin{lstlisting}[style=plainpy]
# B is the BHML companion table; both T and B are recorded in
# Gen13/targets/foundations/tables/canonical_tables.py
from canonical_tables import T, B
print("|T (xor) B|:", int((T != B).sum()))        # 71
\end{lstlisting}

\subsection{Yang--Mills determinant ($70$-rung)}
\label{subsec:ym-det-snippet}

\begin{lstlisting}[style=plainpy]
import numpy as np
from canonical_tables import B
idx = [1, 2, 3, 4, 5, 6, 8, 9]
B_YM = B[np.ix_(idx, idx)]
print("det(B_YM) =", int(round(np.linalg.det(B_YM))))   # 70
print("C(8, 4)   =", int(np.math.comb(8, 4)))           # 70
\end{lstlisting}

\subsection{Verification scripts (bundled)}

The five ladder snippets above are also bundled as standalone scripts
for one-shot reproducibility:

\begin{itemize}[leftmargin=*]
\item \texttt{harmony\_ladder\_disc\_check.py} — verifies the
discriminant claim of §\ref{subsec:disc-snippet} ($\disc = -40896 =
-2^6 \cdot 3^2 \cdot 71$, with the cross-check against the alternate
factorization).
\item \texttt{tsml\_harmony\_count.py} — verifies $\HARM(T) = 73$;
\item \texttt{tsml\_submagma\_9x9.py} — verifies $\HARM(T_{\{1,\ldots,9\}}) = 71$;
\item \texttt{tsml\_bhml\_disagreement.py} — verifies the cell-disagreement count $= 71$;
\item \texttt{bhml\_8\_ym\_det.py} — verifies $\det(B_{\mathrm{YM}}) = 70$.
\end{itemize}

The wrapper module \texttt{harmony\_ladder.py} runs all five scripts
in sequence and emits a $5 \times 3$ verification table
(rung / expected / actual). All five match at integer precision.

\begin{thebibliography}{99}

\bibitem{Bourbaki1968}
N. Bourbaki, \emph{Groupes et alg\`ebres de Lie, Chapitres 4--6},
Hermann, 1968.

\bibitem{DrapalWanless2021}
A. Drápal and I.M. Wanless, \emph{Maximally non-associative
quasigroups}, J.~Combin.~Theory Ser.~A \textbf{184} (2021), 105510.

\bibitem{FultonHarris1991}
W. Fulton and J. Harris, \emph{Representation Theory: A First
Course}, GTM 129, Springer, 1991.

\bibitem{Sanders2026CLAxioms}
B.R. Sanders and M. Gish, \emph{Structural Axioms for the
CL\_TSML Composition Lattice on $\Z/10\Z$}, submitted to
\emph{Algebraic Combinatorics}, 2026 [J16].

\bibitem{Sanders2026LensInvariance}
B.R. Sanders and M. Gish, \emph{TSML 73 Cells / BHML 28 Cells:
Lens-Invariant Cell Counts on the $\Z/10\Z$ Composition Lattice},
submitted to \emph{Experimental Mathematics}, 2026 [J32].

\bibitem{Sanders2026FourCore}
B.R. Sanders and M. Gish, \emph{Joint Closure, Per-Coordinate Fuse
Data, and a Closed-Form Algebraic Attractor of Two Commutative
Binary Operations on $\Z/10\Z$}, submitted to \emph{Algebraic
Combinatorics}, 2026 [J15].

\bibitem{Sanders2026Attractor}
B.R. Sanders and M. Gish, \emph{Galois $D_4$ over LMFDB 4.2.10224.1:
Number-Field Identification of the Four-Core Attractor}, submitted to
\emph{Communications in Algebra}, 2026 [J12].

\bibitem{Sanders2026Wobble}
B.R. Sanders and M. Gish, \emph{On the Prime-Divisibility Pattern of
the Characteristic Polynomial of a $10 \times 10$ Integer Matrix
Arising in a Discrete Magma on $\Z/10\Z$}, submitted to
\emph{Linear Algebra and Its Applications}, 2026 [J19].

\end{thebibliography}

\end{document}
